\begin{pastexam}[label={px:ntust-109-et-ds-06}]{臺科大電子所甲組 109 學年度 第 6 題}
（問答；10 分）

(10\%) What are the average time and worst time to sort $n$ objects using bubble sort, quick sort,
heap sort, merge sort, and radix sort?
\end{pastexam}

\begin{pastsoln}
\begin{center}
\begin{tabular}{lll}
\toprule
排序法 & 平均 & 最壞 \\
\midrule
氣泡排序 & $\Theta(n^2)$ & $\Theta(n^2)$ \\
快速排序 & $\Theta(n\log n)$ & $\Theta(n^2)$ \\
堆積排序 & $\Theta(n\log n)$ & $\Theta(n\log n)$ \\
合併排序 & $\Theta(n\log n)$ & $\Theta(n\log n)$ \\
基數排序 & $\Theta(d(n + r))$ & $\Theta(d(n + r))$ \\
\bottomrule
\end{tabular}
\end{center}
基數排序的 $d$ 是鍵的位數、$r$ 是基數（每一位可能的值有幾種）。 理由（資結書第~14、15~章）：
\begin{itemize}
\item \textbf{氣泡排序}：平均有 $n(n - 1)/4$ 個逆序對，每次相鄰交換只消掉一個，所以平均
$\Theta(n^2)$；最壞（由大到小）要比 $n(n - 1)/2$ 次。 加上「一輪沒有交換就停止」的版本，最好
情形是 $\Theta(n)$，但平均和最壞不變。
\item \textbf{快速排序}：平均每次分割大致均衡，$\Theta(n\log n)$（資結書定理~14.3）；最壞每次都
切成 $0$ 和 $n - 1$（例如固定取第一個當基準，輸入已排序），$T(n) = T(n - 1) + \Theta(n) =
\Theta(n^2)$。
\item \textbf{堆積排序}：建堆積 $O(n)$，再做 $n - 1$ 次取出最大值，每次 $O(\log n)$，不管輸入
如何都是 $O(n\log n)$（資結書第~14.6~節）。
\item \textbf{合併排序}：$T(n) = 2T(n/2) + \Theta(n)$，和輸入無關，$\Theta(n\log n)$（資結書
第~14.4~節）。
\item \textbf{基數排序}（LSD）：每一位用計數排序做一次穩定排序，$\Theta(n + r)$，共 $d$ 位，
$\Theta(d(n + r))$，和輸入順序無關（資結書第~15.2.2~節）。 $d$、$r$ 是常數時是 $\Theta(n)$；
它不靠比較，所以不受 $\Omega(n\log n)$ 下界限制。
\end{itemize}
\end{pastsoln}
